\section{Related Work}
\label{sec:RelatedWork}

\subsection{Molecular Geometry Representation Learning}

Geometric deep learning drives the progress of machine learning force fields (MLFF)~\cite{liu2022spherical,wang2024enhancing}. Equivariant neural networks (ENNs)~\cite{satorras2021n} that respect E(3) symmetries are the dominant paradigm, with methods modeling inter-atomic interactions~\cite{schutt2017schnet,liao2023equiformer}, angular information~\cite{gasteiger2020directional,liu2022spherical,gasteiger2021gemnet}, and higher-order representations~\cite{liao2024equiformerv2,passaro2023reducing,wang2024enhancing}. However, all existing geometry-based methods operate solely on nuclear-level information and do not incorporate electronic structure. Our approach is architecture-agnostic and supplements any geometric model with electronic information.

\subsection{Electron Density Representation Learning}

Existing electron density (ED) representations fall into two categories. Point cloud methods treat all density values as a point set; for example, PointNet~\cite{qi2017pointnet} has been applied to classify the symmetry of inorganic compounds~\cite{kim2024deep}. Voxel methods discretize the density field on a regular grid; 3D CNNs~\cite{liu2015sparse} have been used to predict molecular exchange energy~\cite{gong2023incorporation} and discover host-molecule guests~\cite{parrilla2024electron}, while 3D-UNet~\cite{cciccek20163d} segments reactive sites and classifies substances~\cite{singh2024classification}.

Beyond direct ED representation, several works have explored machine learning approaches to predict or bypass electron density computation. Brockherde~\etal~\cite{brockherde2017bypassing} proposed bypassing the Kohn-Sham equations by directly learning the density-to-energy mapping. Nagai~\etal~\cite{nagai2020completing} used neural networks to learn the exchange-correlation functional from molecular electron densities. Sunshine~\etal~\cite{sunshine2023chemical} predicted chemical properties from GNN-predicted electron densities, demonstrating that learned ED representations carry meaningful chemical information. More recently, Lee~\etal~\cite{lee2025selfsupervised} proposed self-supervised diffusion processes for electron-aware molecular representation learning, and Gao~\etal~\cite{gao2025equivariant} developed equivariant non-local electron density functionals. However, these methods either require ED at inference time, focus on predicting ED itself, or learn ED-aware representations without a cross-modal distillation mechanism. % [Polish] Original: Different from previous methods, we propose a novel multi-view RGB-D image to represent ED and design a representation learning method to extract features, enabling ED knowledge to be transferred to geometry models without requiring ED at inference.
In contrast, we represent ED as multi-view RGB-D images and design a representation learning method to extract features, enabling ED knowledge to be transferred to geometry models without requiring ED at inference.

\subsection{Knowledge Distillation and Privileged Information}

Knowledge distillation (KD)~\cite{hinton2015distilling,gou2021knowledge} transfers knowledge from a teacher model to a student model, typically by matching soft predictions or intermediate representations. Feature-based distillation methods~\cite{romero2015fitnets,zagoruyko2017paying} extend this paradigm by aligning intermediate layer activations. Wang~\etal~\cite{wang2021knowledge} provided a comprehensive review. Recent work has recognized that not all teacher knowledge is equally beneficial. In the selective KD line, SelKD~\cite{park2024selkd} formulates selection via optimal transport, and Shi~\etal~\cite{shi2025selkd} further develop this perspective at scale. Studies on teacher calibration~\cite{guo2017calibration,das2023understanding} and uncertainty-aware teaching~\cite{lin2023teaching,guo2023selective} have shown that accounting for the teacher's confidence can significantly improve outcomes.

Cross-modal KD is increasingly studied~\cite{huo2024c2kd}, showing that bridging heterogeneous feature spaces requires explicit alignment. Our framework distills between fundamentally different representations---ED images and atomic graphs---where the representational and symmetry gap poses distinct challenges.

\textbf{Positioning of EDG++.} We distinguish four sources of unreliable supervision: (1) \emph{sample difficulty}~\cite{bengio2009curriculum,kumar2010self}; (2) \emph{label noise}~\cite{han2018co,song2022learning}; (3) \emph{teacher uncertainty}~\cite{park2024selkd,shi2025selkd}; (4) \emph{cross-modal alignment quality}---where the cross-modal mapping itself is imperfect for certain samples. EDG++ operates in category (4): task labels are clean, but the quality of the cross-modal prediction varies due to information loss in the 3D-to-2D projection. This motivates a reliability estimator that predicts the teacher's reconstruction error, rather than relying on the teacher's own uncertainty. Our work builds upon the LUPI paradigm~\cite{vapnik2009new,lopez2016unifying}, with electron density as the privileged modality. Additional discussion of selective learning methods is provided in Appendix~\ref{sec:appendix_selective_learning}.
